\documentclass[../DoAn.tex]{subfiles}
\begin{document}

% ═══════════════════════════════════════════════════════════════════════════════
\section{Blockchain Ethereum, Mô hình Tài khoản và ERC-4337}
% ═══════════════════════════════════════════════════════════════════════════════

\subsection{Mô hình Tài khoản Ethereum}

Ethereum quản lý trạng thái toàn cục theo \textbf{mô hình tài khoản (Account Model)}, trong đó mỗi tài khoản được định danh bởi một địa chỉ 20-byte và lưu trữ bộ bốn trạng thái: số dư ETH (\texttt{balance}), bộ đếm giao dịch (\texttt{nonce}), mã bytecode (\texttt{codeHash}), và gốc cây trạng thái lưu trữ (\texttt{storageRoot})~\cite{ethereum_whitepaper}. Tồn tại hai loại tài khoản chính:

\begin{itemize}
    \item \textbf{Externally Owned Accounts (EOA):} Được kiểm soát bởi người dùng thông qua cặp khóa \((sk, pk)\) trên đường cong \texttt{secp256k1}, trong đó \(pk = sk \cdot G\). Mọi giao dịch phát sinh từ EOA phải được ký bởi khóa riêng \(sk\), và phí Gas được thanh toán trực tiếp bằng ETH từ tài khoản đó. Đây là loại tài khoản duy nhất có thể khởi tạo giao dịch gốc trên mạng.

    \item \textbf{Contract Accounts:} Là các tài khoản được điều khiển bởi mã bytecode EVM, không có khóa riêng tương ứng. Chúng chỉ thực thi logic khi nhận được lời gọi từ giao dịch EOA hoặc từ một hợp đồng khác. Địa chỉ hợp đồng được xác định tại thời điểm triển khai, hoặc tính toán trước qua cơ chế \texttt{CREATE2} (xem Mục~\ref{subsec:create2}).
\end{itemize}

Do đặc tính minh bạch của blockchain công khai, toàn bộ lịch sử giao dịch, luồng tiền và số dư tài khoản đều có thể bị quan sát bởi bất kỳ bên nào. Điều này tạo ra rủi ro nghiêm trọng về quyền riêng tư khi người dùng muốn thực hiện giao dịch ẩn danh.

\subsection{Account Abstraction theo ERC-4337}
\label{subsec:erc4337}

Để khắc phục sự cứng nhắc của EOA truyền thống và mở rộng khả năng kiểm soát tài khoản, hệ thống áp dụng mô hình \textbf{Account Abstraction (AA)} theo tiêu chuẩn ERC-4337~\cite{erc4337}. Khác biệt cốt lõi so với EOA là: logic xác thực giao dịch được chuyển hoàn toàn vào bên trong Smart Contract chứ không phụ thuộc vào chữ ký ECDSA cố định. Kiến trúc này hoạt động thông qua ba thành phần chính:

\begin{enumerate}
    \item \textbf{UserOperation:} Thay vì giao dịch thông thường, người dùng tạo ra một cấu trúc dữ liệu \texttt{UserOperation} bao gồm:
    \begin{itemize}
        \item \texttt{sender}: địa chỉ Smart Account gửi yêu cầu.
        \item \texttt{nonce}: bộ đếm chống phát lại.
        \item \texttt{initCode}: bytecode dùng để triển khai tài khoản lần đầu (nếu chưa tồn tại).
        \item \texttt{callData}: dữ liệu lời gọi hàm cần thực thi.
        \item \texttt{paymasterAndData}: địa chỉ Paymaster và dữ liệu tùy chọn để tài trợ Gas.
        \item \texttt{signature}: dữ liệu xác thực tùy chỉnh -- trong hệ thống này là bằng chứng ZK (xem Mục~\ref{subsec:zk_spend}).
    \end{itemize}

    \item \textbf{Bundler và EntryPoint:} \texttt{UserOperation} được gửi tới một \textbf{Bundler} ngoài chuỗi. Bundler tổng hợp nhiều \texttt{UserOperation} từ các tài khoản khác nhau và gọi hàm \texttt{handleOps()} trên hợp đồng \texttt{EntryPoint} tiêu chuẩn. \texttt{EntryPoint} lần lượt gọi \texttt{validateUserOp()} trên từng Smart Account để xác thực, sau đó thực thi lời gọi.

    \item \textbf{Paymaster hỗ trợ phí Gas:} Hợp đồng \texttt{StealthPaymaster} (xem Mục~\ref{subsec:paymaster}) cho phép tài trợ phí Gas thay cho tài khoản người dùng. Điều này giải quyết bài toán ``Cold Start'': một Stealth Account mới tạo chưa có ETH nhưng vẫn có thể thực hiện giao dịch đầu tiên nhờ Paymaster bảo lãnh phí Gas.
\end{enumerate}

Hình~\ref{fig:erc4337_flow} mô tả luồng xử lý \texttt{UserOperation} trong kiến trúc ERC-4337.

\begin{figure}[h]
\centering
\includegraphics[width=0.85\textwidth]{Hinhve/erc4337-transaction-flow-800x450.png}
\caption{Luồng xử lý UserOperation trong kiến trúc ERC-4337~\cite{erc4337}}
\label{fig:erc4337_flow}
\end{figure}

\subsection{StealthPaymaster và giải pháp Cold Start}
\label{subsec:paymaster}

Trong hệ thống Stealth Wallet, mỗi địa chỉ stealth mới được tạo ra là một tài khoản chưa có ETH. Nếu người dùng muốn chi tiêu từ tài khoản này, họ cần có ETH để trả Gas -- nhưng để có ETH, họ cần thực hiện giao dịch trước. Đây là vòng luẩn quẩn ``Cold Start''.

Hợp đồng \texttt{StealthPaymaster} giải quyết vấn đề này bằng cách:
\begin{enumerate}
    \item Triển khai hàm \texttt{validatePaymasterUserOp()} để \texttt{EntryPoint} gọi khi xử lý \texttt{UserOperation}. Hàm này luôn trả về \texttt{validationData = 0} (hợp lệ) cho \textit{mọi} \texttt{UserOperation}, tức là bảo lãnh không hạn chế.
    \item Nạp ETH vào \texttt{EntryPoint} thông qua hàm \texttt{deposit()}, duy trì số dư đủ để tài trợ phí Gas cho các giao dịch stealth.
\end{enumerate}

Thiết kế này thích hợp cho môi trường nghiên cứu trên Sepolia Testnet. Trong môi trường sản xuất, cần bổ sung cơ chế kiểm soát truy cập và giới hạn tài trợ.

% ═══════════════════════════════════════════════════════════════════════════════
\section{Cơ sở Toán học và Mật mã học Đường cong Elliptic}
% ═══════════════════════════════════════════════════════════════════════════════

\subsection{Bài toán Logarit Rời rạc trên Đường cong Elliptic (ECDLP)}

Nền tảng bảo mật của hệ thống dựa trên độ khó của \textbf{Bài toán Logarit Rời rạc trên Đường cong Elliptic (ECDLP)}~\cite{ecc}: cho trước điểm sinh \(G\) và điểm công khai \(P = k \cdot G\) trên đường cong, về mặt tính toán không thể xác định giá trị nguyên \(k\) trong thời gian đa thức. Tính chất nhất chiều này là nền tảng cho giao thức ECDH và hệ thống chữ ký ECDSA.

\subsection{Hai đường cong được sử dụng}

Hệ thống sử dụng đồng thời hai đường cong elliptic với vai trò khác nhau:

\begin{enumerate}
    \item \textbf{secp256k1} -- Đường cong chuẩn của Ethereum~\cite{erc5564}:
    \begin{itemize}
        \item Phương trình: \(y^2 = x^3 + 7 \pmod{p}\), với \(p = 2^{256} - 2^{32} - 977\).
        \item Bậc nhóm: \(n \approx 1.16 \times 10^{77}\).
        \item Được dùng cho: tạo khóa xem/chi tiêu, giao thức ECDH để dẫn xuất Shared Secret, và xác minh quyền sở hữu địa chỉ Ethereum (EIP-55).
    \end{itemize}

    \item \textbf{BN254 (Baby-Jubjub / BN-128)} -- Đường cong thân thiện với pairing~\cite{groth16}:
    \begin{itemize}
        \item Được dùng trong mạch Circom và hàm băm Poseidon để xây dựng Merkle Tree.
        \item Trường hữu hạn: \(\mathbb{F}_p\) với \(p = 21888242871839275222246405745257275088548364400416034343698204186575808495617\).
        \item Tương thích với hệ thống chứng minh Groth16 trên EVM thông qua precompile \texttt{ecPairing} (EIP-197).
    \end{itemize}
\end{enumerate}

Sự tách biệt hai đường cong là cần thiết vì: secp256k1 không có cấu trúc pairing-friendly nên không thể dùng trực tiếp trong zk-SNARK Groth16, trong khi BN254 không tương thích với chữ ký Ethereum.

\subsection{Giao thức Trao đổi Khóa Diffie-Hellman trên Đường cong Elliptic (ECDH)}
\label{subsec:ecdh}

Giao thức ECDH cho phép hai bên thiết lập một shared secret chung qua kênh truyền công khai mà không cần trao đổi bí mật trực tiếp~\cite{ecc}:

\begin{enumerate}
    \item Người nhận tạo cặp khóa xem: khóa riêng tư \(v \xleftarrow{\$} \mathbb{Z}_n^*\) và khóa công khai \(V = v \cdot G\).
    \item Người gửi sinh số ngẫu nhiên tạm thời \(r \xleftarrow{\$} \mathbb{Z}_n^*\), tính:
    \[
        S_{\text{shared}} = r \cdot V = r \cdot v \cdot G
    \]
    \item Người nhận tái tạo shared secret bằng khóa riêng tư của mình:
    \[
        S'_{\text{shared}} = v \cdot R = v \cdot r \cdot G = S_{\text{shared}}
    \]
    nơi \(R = r \cdot G\) là khóa công khai tạm thời được người gửi công bố.
\end{enumerate}

Tính đúng đắn: \(S_{\text{shared}} = S'_{\text{shared}}\) vì nhân vô hướng trên nhóm giao hoán. Tính bảo mật: \(S_{\text{shared}}\) không thể tái tạo bởi bên thứ ba chỉ biết \(R\) và \(V\) mà không biết \(r\) hoặc \(v\), dựa trên giả thiết ECDLP.

% ═══════════════════════════════════════════════════════════════════════════════
\section{Giao thức Stealth Address (ERC-5564)}
% ═══════════════════════════════════════════════════════════════════════════════

\subsection{Tổng quan Giao thức}
\label{subsec:stealth_protocol}

Giao thức Stealth Address theo ERC-5564~\cite{erc5564} cho phép người gửi chuyển tài sản đến một địa chỉ một lần (one-time address) mà chỉ người nhận mới có thể xác định và kiểm soát, đảm bảo không có liên kết quan sát được trên chuỗi giữa danh tính người nhận và địa chỉ nhận tiền.

Trong kiến trúc đề xuất, giao thức ERC-5564 được mở rộng để tích hợp với hệ thống Merkle Tree và ZK Proof: quyền chi tiêu không phụ thuộc vào ECDSA thuần túy mà được xác minh thông qua shared secret kết hợp với dữ liệu định danh trong Merkle Tree. Điều này cho phép tách biệt hoàn toàn ``quyền xem'' (\textit{view authority}) khỏi ``quyền chi tiêu'' (\textit{spend authority}).

\subsection{Tạo Meta-address}
\label{subsec:meta_address}

Người nhận (Bob) khởi tạo \textbf{Stealth Meta-address} như sau:

\begin{enumerate}
    \item Sinh cặp khóa xem (\textit{view key pair}):
    \[
        v \xleftarrow{\$} \mathbb{Z}_n^*, \quad V = v \cdot G \quad \text{(trên secp256k1)}
    \]
    
    \item Sinh cặp khóa chi tiêu (\textit{spend key pair}):
    \[
        x \xleftarrow{\$} \mathbb{F}_p, \quad \text{Leaf} = \text{Poseidon}(x) \quad \text{(trên BN254)}
    \]
    
    \item Đăng ký lá vào Merkle Tree thông qua Relayer Server. Sau khi đăng ký, hệ thống trả về chỉ mục \(i\) của lá trong cây, và tính:
    \[
        \text{IndexHash} = \text{Poseidon}(i, 0)
    \]

    \item Công bố Meta-address bao gồm: \(V\) (khóa xem công khai), \(\text{spendPub}\) (khóa chi tiêu công khai) và \(\text{IndexHash}\). Ba giá trị này cũng được liên kết với tên miền ENS của người dùng (xem Mục~\ref{subsec:ens}).
\end{enumerate}

\subsection{Quá trình gửi (Sender Side)}
\label{subsec:sending}

Người gửi (Alice) thực hiện các bước sau để gửi tài sản riêng tư đến Bob:

\begin{enumerate}
    \item Truy xuất Meta-address của Bob: \((V_{\text{Bob}}, \text{spendPub}_{\text{Bob}}, \text{IndexHash}_{\text{Bob}})\).

    \item Sinh số ngẫu nhiên tạm thời và tính Shared Secret:
    \[
        r \xleftarrow{\$} \mathbb{Z}_n^*, \quad R = r \cdot G, \quad S_{\text{shared}} = r \cdot V_{\text{Bob}}
    \]

    \item Tính \texttt{sharedSecretHash}:
    \[
        h = \text{keccak256}(S_{\text{shared}}.x \,\|\, S_{\text{shared}}.y)
    \]

    \item Tính \textbf{Index Commitment} -- giá trị liên kết địa chỉ stealth với lá Merkle:
    \[
        C_{\text{index}} = \text{Poseidon}\bigl(\text{IndexHash}_{\text{Bob}} + h\bigr)
    \]

    \item Tính địa chỉ Stealth Account bằng \texttt{CREATE2} (xem Mục~\ref{subsec:create2}):
    \[
        A_{\text{stealth}} = \text{Factory.getAddress}(C_{\text{index}})
    \]

    \item Chuyển tài sản đến \(A_{\text{stealth}}\) và phát sự kiện \texttt{Announcement} chứa:
    \begin{itemize}
        \item \(R = r \cdot G\): khóa công khai tạm thời để Bob quét.
        \item \(\text{schemeId}\): định danh giao thức (ERC-5564 scheme 1).
        \item \texttt{metadata}: thông tin mã hóa về loại tài sản.
    \end{itemize}
\end{enumerate}

\subsection{Quá trình quét (Receiver Side)}
\label{subsec:scanning}

Bob liên tục lắng nghe sự kiện \texttt{Announcement} từ hợp đồng \texttt{ERC5564Announcer}. Với mỗi sự kiện, Bob:

\begin{enumerate}
    \item Tái tạo Shared Secret:
    \[
        S'_{\text{shared}} = v_{\text{Bob}} \cdot R = v_{\text{Bob}} \cdot r \cdot G = S_{\text{shared}}
    \]

    \item Tính \texttt{sharedSecretHash} và \texttt{indexCommitment} tương tự như phía người gửi.

    \item Tính địa chỉ Stealth Account kỳ vọng \(A'_{\text{stealth}}\) bằng \texttt{CREATE2}.

    \item So sánh \(A'_{\text{stealth}}\) với địa chỉ trong sự kiện. Nếu khớp, Bob xác định tài khoản này thuộc về mình và ghi nhớ \(\texttt{sharedSecretHash}\) phục vụ cho việc tạo bằng chứng ZK sau này.
\end{enumerate}

% ═══════════════════════════════════════════════════════════════════════════════
\section{Kiến trúc CREATE2 và Triển khai Địa chỉ Định trước}
% ═══════════════════════════════════════════════════════════════════════════════

\subsection{Cơ chế CREATE2}
\label{subsec:create2}

Opcode \texttt{CREATE2} (EIP-1014)~\cite{eip1014} cho phép xác định trước địa chỉ của một hợp đồng sẽ được triển khai trong tương lai, theo công thức:

\[
    A_{\text{stealth}} = \text{keccak256}\!\left(
        \texttt{0xff} \;\|\; \text{factory} \;\|\; \text{salt} \;\|\; \text{keccak256(bytecode)}
    \right)[12:]
\]

trong đó:
\begin{itemize}
    \item \(\text{factory}\): địa chỉ hợp đồng \texttt{StealthAccountFactory}.
    \item \(\text{salt} = C_{\text{index}}\): giá trị Index Commitment 32-byte.
    \item \(\text{bytecode}\): creation bytecode của \texttt{StealthAccount} được encode kèm các tham số khởi tạo \((C_{\text{index}}, \text{treeManager}, \text{poseidonHasher}, \text{verifier})\).
\end{itemize}

\textbf{Tính đơn định (deterministic):} Địa chỉ \(A_{\text{stealth}}\) được xác định hoàn toàn bởi \(C_{\text{index}}\) và địa chỉ factory. Do đó, cả người gửi lẫn người nhận đều có thể độc lập tính toán địa chỉ này mà không cần trao đổi thêm thông tin.

\textbf{Tính lười triển khai (lazy deployment):} Tài sản có thể được gửi đến \(A_{\text{stealth}}\) trước khi hợp đồng được triển khai (trạng thái ``counterfactual''). Hợp đồng chỉ được triển khai thực sự khi và chỉ khi người dùng muốn chi tiêu.

\subsection{Liên kết với Merkle Tree}
\label{subsec:merkle_link}

Giá trị salt được sử dụng trong \texttt{CREATE2} không phải là một giá trị tùy ý mà được dẫn xuất có cấu trúc:
\[
    \text{salt} = C_{\text{index}} = \text{Poseidon}(\text{IndexHash}_{\text{Bob}} + \text{sharedSecretHash})
\]

Cấu trúc này đảm bảo ba tính chất:
\begin{enumerate}
    \item \textbf{Ràng buộc:} Địa chỉ Stealth Account được liên kết toán học với lá tại vị trí \(i\) trong Merkle Tree, ngăn chặn việc giả mạo quyền chi tiêu.
    \item \textbf{Bảo mật danh tính:} Vị trí lá \(i\) và khóa chi tiêu \(x\) không được công khai; chỉ \(C_{\text{index}}\) (một cam kết ẩn) mới xuất hiện trên chuỗi.
    \item \textbf{Tính nhất quán:} Cùng một giá trị \(C_{\text{index}}\) được sử dụng làm tín hiệu công khai trong mạch ZK, đảm bảo chứng minh và hợp đồng đồng thuận về danh tính.
\end{enumerate}

% ═══════════════════════════════════════════════════════════════════════════════
\section{Hàm Băm Poseidon và Cây Merkle}
% ═══════════════════════════════════════════════════════════════════════════════

\subsection{Hàm Băm Poseidon}
\label{subsec:poseidon}

Poseidon là hàm băm mật mã được thiết kế đặc biệt để tối ưu hóa trong các hệ thống chứng minh số học như zk-SNARK~\cite{poseidon}. So với các hàm băm truyền thống (SHA-256, Keccak), Poseidon hoạt động trực tiếp trên các phần tử của trường hữu hạn \(\mathbb{F}_p\), cho phép số ràng buộc R1CS (Rank-1 Constraint System) thấp hơn nhiều bậc.

Hệ thống sử dụng Poseidon với hai cấu hình:
\begin{itemize}
    \item \(\text{Poseidon}(1)\): tính băm của một phần tử -- dùng để tính \(\text{Leaf} = \text{Poseidon}(x)\) từ khóa chi tiêu.
    \item \(\text{Poseidon}(2)\): tính băm của hai phần tử -- dùng để kết hợp các nút trong Merkle Tree và tính chỉ mục cam kết.
\end{itemize}

\subsection{Cấu trúc Cây Merkle Poseidon}
\label{subsec:merkle_tree}

Hệ thống duy trì một \textbf{Incremental Merkle Tree} độ sâu \(d = 20\) (tối đa \(2^{20} \approx 10^6\) lá), sử dụng Poseidon làm hàm nội tại. Cây được quản lý ngoài chuỗi bởi Relayer Server và chỉ lưu trữ giá trị gốc (\texttt{root}) trên chuỗi.

Cấu trúc mỗi lá:
\[
    \text{Leaf}_i = \text{Poseidon}(x_i)
\]
nơi \(x_i\) là khóa chi tiêu riêng của người dùng \(i\). Giá trị lá được giữ bí mật hoàn toàn; chỉ \(\text{IndexHash} = \text{Poseidon}(i, 0)\) mới được công khai.

Gốc Merkle:
\[
    R_{\text{tree}} = H_{\text{Poseidon}}(\cdots H_{\text{Poseidon}}(\text{Leaf}_0, \text{Leaf}_1), \ldots)
\]

được cập nhật trên chuỗi mỗi khi một lá mới được thêm vào hoặc cập nhật (trong quá trình phục hồi).

% ═══════════════════════════════════════════════════════════════════════════════
\section{Zero-Knowledge Proof: Groth16 và Mạch Circom}
% ═══════════════════════════════════════════════════════════════════════════════

\subsection{Tổng quan Groth16}
\label{subsec:groth16}

Groth16~\cite{groth16} là hệ thống chứng minh SNARK (\textit{Succinct Non-interactive ARgument of Knowledge}) dựa trên cặp đôi đường cong elliptic (bilinear pairing). Với một chương trình tính toán được mã hóa dưới dạng mạch R1CS, Groth16 cho phép:
\begin{itemize}
    \item \textbf{Prover:} Tạo bằng chứng \(\pi = (A, B, C)\) kích thước không đổi (3 điểm trên đường cong, tổng 192 byte) trong thời gian \(O(n \log n)\), nơi \(n\) là số ràng buộc.
    \item \textbf{Verifier:} Xác minh bằng chứng \(\pi\) với tín hiệu công khai \(\vec{x}\) bằng 3 phép tính pairing trong \(O(1)\) -- phù hợp để thực hiện trên EVM thông qua precompile \texttt{bn256Pairing}.
\end{itemize}

Đặc điểm quan trọng của Groth16:
\begin{itemize}
    \item \textbf{Zero-knowledge:} Bằng chứng không tiết lộ bất kỳ thông tin nào về đầu vào riêng tư.
    \item \textbf{Soundness:} Không thể tạo bằng chứng hợp lệ cho một câu lệnh sai (với xác suất bỏ qua).
    \item \textbf{Trusted Setup:} Yêu cầu một lần khởi tạo tham số (Powers-of-Tau ceremony) cho mỗi mạch.
\end{itemize}

\subsection{Mạch ZK cho Chi tiêu (\texttt{stealth\_spend})}
\label{subsec:zk_spend}

Mạch \texttt{StealthAccountZK} (độ sâu \(d = 20\)) được triển khai trong Circom với cấu trúc như sau:

\textbf{Tín hiệu công khai (public inputs):}
\begin{itemize}
    \item \(\texttt{root} \in \mathbb{F}_p\): Gốc Merkle Tree hiện tại trên chuỗi.
    \item \(\texttt{indexCommitment} \in \mathbb{F}_p\): Cam kết định danh tài khoản \(C_{\text{index}}\) được lưu trong hợp đồng.
\end{itemize}

\textbf{Tín hiệu riêng tư (private inputs):}
\begin{itemize}
    \item \(x\): Khóa chi tiêu riêng tư (\texttt{spendPriv}).
    \item \(\texttt{sharedSecretHash}\): Băm của Shared Secret từ giao dịch nhận.
    \item \(\texttt{merkleProof}[20]\): Đường chứng minh Merkle (20 nút anh em).
    \item \(\texttt{pathIndices}\): Chỉ mục vị trí lá trong cây (encode nhị phân).
\end{itemize}

\textbf{Các ràng buộc (constraints):}

\begin{enumerate}
    \item \textbf{Ràng buộc tính lá:}
    \[
        k = \text{Poseidon}(x)
    \]
    Đảm bảo người dùng biết khóa chi tiêu \(x\) tương ứng với lá \(k\) trong cây.

    \item \textbf{Ràng buộc cam kết chỉ mục:}
    \[
        \text{indexHash} = \text{Poseidon}(\texttt{pathIndices},\; 0)
    \]
    \[
        \texttt{indexCommitment} = \text{Poseidon}(\text{indexHash} + \texttt{sharedSecretHash})
    \]
    Đảm bảo rằng \texttt{indexCommitment} public trên chuỗi khớp với vị trí lá thực sự và shared secret của giao dịch nhận.

    \item \textbf{Ràng buộc Merkle Inclusion:}
    \[
        \text{MerkleTreeChecker}(k,\; \texttt{root},\; \texttt{merkleProof},\; \texttt{pathIndices}) = \text{true}
    \]
    Xác minh rằng lá \(k\) tồn tại hợp lệ trong Merkle Tree hiện tại có gốc \texttt{root}.
\end{enumerate}

\textbf{Quá trình xác minh trên chuỗi:} Hợp đồng \texttt{StealthSpendVerifier} nhận bằng chứng \(\pi = (A, B, C)\) và hai tín hiệu công khai \((\texttt{root}, \texttt{indexCommitment})\), thực hiện xác minh pairing. Nếu thành công, \texttt{StealthAccount.execute()} cho phép thực hiện giao dịch.

\subsection{Mạch ZK cho Cập nhật Cây (\texttt{smt\_update})}
\label{subsec:zk_smt}

Mạch \texttt{SMTUpdate} chứng minh rằng một chuỗi cập nhật lá trong Merkle Tree là hợp lệ mà không cần thực hiện toàn bộ tính toán cây trên EVM.

\textbf{Tín hiệu công khai:}
\begin{itemize}
    \item \(\texttt{oldRoot}\): Gốc Merkle trước khi cập nhật.
    \item \(\texttt{newRoot}\): Gốc Merkle sau khi cập nhật.
    \item \(\texttt{index}\): Vị trí lá được cập nhật.
    \item \(\texttt{newLeaf}\): Giá trị lá mới.
\end{itemize}

\textbf{Tín hiệu riêng tư:}
\begin{itemize}
    \item \(\texttt{oldLeaf}\): Giá trị lá cũ tại vị trí \texttt{index}.
    \item \(\texttt{siblings}[20]\): Các nút anh em trên đường Merkle (không thay đổi).
\end{itemize}

\textbf{Các ràng buộc:}

Mạch xác minh đồng thời hai ràng buộc Merkle sử dụng cùng một tập \(\texttt{siblings}\):
\[
    f_{\text{merkle}}(\texttt{oldLeaf},\; \texttt{index},\; \texttt{siblings}) = \texttt{oldRoot}
\]
\[
    f_{\text{merkle}}(\texttt{newLeaf},\; \texttt{index},\; \texttt{siblings}) = \texttt{newRoot}
\]

Tính đúng đắn: Nếu cả hai ràng buộc đồng thời thỏa mãn với cùng \(\texttt{siblings}\), chứng tỏ rằng \(\texttt{newRoot}\) thu được bằng cách thay thế đúng lá tại \(\texttt{index}\) từ \(\texttt{oldRoot}\), và mọi nút còn lại không bị thay đổi. Điều này giúp giảm chi phí xác minh trên chuỗi từ \(O(d)\) phép băm Poseidon xuống \(O(1)\) phép kiểm tra pairing.

% ═══════════════════════════════════════════════════════════════════════════════
\section{Phục hồi Xã hội với Ủy ban Guardian}
% ═══════════════════════════════════════════════════════════════════════════════

\subsection{Vấn đề Phục hồi Tài khoản Stealth}
\label{subsec:recovery_problem}

Tài khoản Stealth được bảo vệ bởi khóa chi tiêu riêng \(x\), không có cơ chế khôi phục mật khẩu truyền thống. Nếu người dùng mất \(x\), toàn bộ tài sản trong tài khoản trở nên không thể truy cập. Đây là vấn đề cố hữu của mô hình ``self-custody''.

Đồng thời, cơ chế phục hồi không được tiết lộ danh tính guardian để tránh tạo điểm tấn công xã hội (social engineering). Hệ thống giải quyết điều này bằng mô hình \textbf{Social Recovery với Cam kết Chỉ mục (Index Commitment)}.

\subsection{Kiến trúc Phục hồi Xã hội}
\label{subsec:social_recovery}

Hợp đồng \texttt{SocialRecoveryModule} được triển khai riêng cho mỗi Stealth Account, liên kết với vị trí lá \(\texttt{mappedIndex}\) trong Merkle Tree. Cơ chế hoạt động:

\begin{enumerate}
    \item \textbf{Thiết lập Guardian:} Chủ sở hữu cung cấp danh sách địa chỉ guardian và ngưỡng phê duyệt \(t\) (threshold). Danh sách này được lưu dưới dạng mapping \texttt{isGuardian[address] => bool} với tổng số \(n\) guardian, yêu cầu \(t \leq n\).

    \item \textbf{Đề xuất Phục hồi:} Bất kỳ guardian nào cũng có thể gọi \texttt{proposeRecovery(newRoot, newLeaf)} để đề xuất cập nhật lá tại \texttt{mappedIndex}:
    \begin{itemize}
        \item \(\texttt{newLeaf} = \text{Poseidon}(x_{\text{mới}})\): cam kết khóa chi tiêu mới.
        \item \(\texttt{newRoot}\): gốc Merkle mới sau khi thay thế lá.
    \end{itemize}
    Guardian đề xuất cũng tự động phê duyệt yêu cầu.

    \item \textbf{Bỏ phiếu Phê duyệt:} Các guardian còn lại gọi \texttt{approveRecovery(reqId)} để phê duyệt. Hệ thống đảm bảo mỗi guardian chỉ bỏ phiếu một lần mỗi yêu cầu.

    \item \textbf{Thực thi Phục hồi:} Khi số phê duyệt \(\geq t\), bất kỳ guardian nào gọi \texttt{executeRecovery(reqId, auth)} với bằng chứng ZK từ mạch \texttt{SMTUpdate}:
    \[
        \texttt{treeManager.updateRoot(newRoot, newLeaf, mappedIndex, auth)}
    \]
    Hợp đồng \texttt{TreeManager} xác minh bằng chứng ZK trước khi cập nhật gốc trên chuỗi.
\end{enumerate}

\textbf{Tính bảo mật danh tính:} Danh sách guardian được lưu trực tiếp theo địa chỉ (không phải cam kết hash), do đó quyền riêng tư của guardian phụ thuộc vào địa chỉ Ethereum không bị liên kết với danh tính thực. Trong các phiên bản nâng cao, có thể thay thế bằng cam kết \(\text{keccak256}(\text{guardian} \,\|\, \text{secret})\) và Merkle Tree riêng để đảm bảo ẩn danh hoàn toàn.

% ═══════════════════════════════════════════════════════════════════════════════
\section{Tích hợp ENS và Định danh Người dùng}
% ═══════════════════════════════════════════════════════════════════════════════
\label{subsec:ens}

\subsection{Ethereum Name Service (ENS)}

ENS~\cite{ens} là hệ thống tên miền phi tập trung trên Ethereum, cho phép ánh xạ tên miền dạng \texttt{name.eth} đến bất kỳ dữ liệu blockchain nào thông qua các bản ghi resolver. Hệ thống hỗ trợ đăng ký tên miền ENS trực tiếp từ giao diện người dùng trên Sepolia Testnet thông qua hợp đồng \texttt{ETHRegistrarController}.

\subsection{Liên kết Meta-address với ENS}

Sau khi đăng ký tên miền, người dùng thiết lập ba bản ghi văn bản (text records) trên ENS Resolver:
\begin{itemize}
    \item \texttt{stealth.scanPub}: Khóa xem công khai \(V = v \cdot G\) (hex-encoded).
    \item \texttt{stealth.spendPub}: Khóa chi tiêu công khai \(\text{spendPub}\) (hex-encoded).
    \item \texttt{stealth.indexHash}: Giá trị \(\text{IndexHash} = \text{Poseidon}(i, 0)\) (hex-encoded).
\end{itemize}

Người gửi có thể tra cứu ba giá trị này chỉ từ tên miền ENS của người nhận, thay vì yêu cầu trao đổi trực tiếp khóa công khai. Thiết kế này tuân thủ tiêu chuẩn ERC-5564 về khả năng tương tác và tính phổ cập của stealth address.

\subsection{Quy trình Đăng ký ENS Tự động}

Quy trình đăng ký ENS trong hệ thống tuân theo giao thức hai bước bắt buộc của ENS để ngăn chặn front-running:

\begin{enumerate}
    \item \textbf{Bước Commit:} Tính toán và gửi cam kết \(\texttt{commitment} = H(\text{name}, \text{owner}, \text{secret})\) lên chuỗi. Người dùng phải đợi tối thiểu 60 giây sau bước này.
    \item \textbf{Bước Register:} Gửi giao dịch đăng ký kèm \(\text{secret}\) và thời gian đăng ký (tính bằng giây, tương ứng với số năm sở hữu) và giá trị ETH trả phí.
\end{enumerate}

Sau khi đăng ký thành công, hệ thống tự động gọi \texttt{setText()} trên Resolver để ghi ba bản ghi stealth meta-address trong cùng một giao dịch.

\end{document}